% 自编教学补充；阅读来源见本章 README。
\begin{frame}[fragile]{6.2A enumerate：编号和内容一起取}
\small 给签到名单编号时，不需要自己维护计数器。enumerate 的 start=1 只改变显示的编号，不改变列表索引。
\medskip
\begin{lstlisting}
attendees = ["小林", "小周", "小陈"]
for number, person in enumerate(attendees, start=1):
    print(number, person)
\end{lstlisting}
\end{frame}

\begin{frame}[fragile]{6.2B zip：把姓名与成绩配对}
\small zip 按位置配对；默认到最短序列就停止。配对前检查长度，避免悄悄漏掉学生。这与 pandas 按标签对齐不同。
\medskip
\begin{lstlisting}
students = ["小林", "小周", "小陈"]
scores = [82, 95, 68]
print(len(students), len(scores))
for student, score in zip(students, scores):
    print(f"{student}: {score}分")
score_map = dict(zip(students, scores))
print(score_map)
\end{lstlisting}
\end{frame}
\begin{frame}[fragile]{课堂练习：6.2B zip：把姓名与成绩配对}
输出分数至少为 80 的学生姓名。删去 scores 的最后一个元素，观察配对数量。
\medskip
\begin{block}{检查思路}
先写出预期结果，再运行。参考答案见配套 Notebook 的折叠单元格。
\end{block}
\end{frame}

\begin{frame}[fragile]{6.2C 列表推导式：从已有循环缩写}
\small 先理解循环，再使用推导式。读法是“对每个 score，若达到 80，就保留 score”。复杂流程仍分步写。
\medskip
\begin{lstlisting}
scores = [82, 95, 68]
passed = []
for score in scores:
    if score >= 80:
        passed.append(score)
short_form = [score for score in scores if score >= 80]
print(passed, short_form)
\end{lstlisting}
\end{frame}
